The commutator identity is the bilinear computation used already in \cite[Prop.~6.1]{ButlerT7} (there for the currents; here for general index pairs): $[\psi_i^\dagger\psi_j, \psi_k^\dagger\psi_l] = \delta_{jk}\psi_i^\dagger\psi_l - \delta_{li}\psi_k^\dagger\psi_j$ from the CAR (\claimref{std-algebra.witt-basis-car}). Each $E_{ij}$ maps the one-particle subspace $\operatorname{span}\{\psi_k^\dagger|0\rangle\}$ to itself, acting as $e_{ij}$ ($E_{ij}\psi_k^\dagger|0\rangle = \delta_{jk}\psi_i^\dagger|0\rangle$); the matrix units generate $\Mat_n$, so the restriction map is onto. The doubling is what supplies the complex structure: each real mode with its adjoined partner supplies the complex $\psi_i, \psi_i^\dagger$ from which the gauge-invariant $E_{ij}$ are built; real Clifford directions alone give the orthogonal bilinears, and it is the pairing that produces $\mathfrak{gl}(n)$. The script checks the commutator table and the one-particle action exactly on a five-mode instance.
